\documentclass[11pt]{article}
\usepackage[T1]{fontenc}
\usepackage{lmodern,microtype,amsmath,amssymb,amsthm,booktabs}
\microtypesetup{expansion=false}
\usepackage[margin=1.05in,headheight=14pt]{geometry}
\usepackage{xcolor}
\definecolor{linknavy}{RGB}{25,55,125}
\usepackage[colorlinks=true,linkcolor=linknavy,citecolor=linknavy,urlcolor=linknavy]{hyperref}
\usepackage{orcidlink,fancyhdr}
\hypersetup{
  pdftitle={A sharp three-point kernel bound and improved proportions of zeta zeros},
  pdfauthor={Felipe Santibanez-Leal},
  pdfsubject={A sharp auxiliary ratio in Wang's three-point kernel argument and its consequences for simple critical and distinct zeta zeros},
  pdfkeywords={Riemann zeta function, simple zeros, distinct zeros, critical line, three-point kernel, spectral defect, interval arithmetic}
}
\newtheorem{theorem}{Theorem}[section]
\newtheorem{proposition}[theorem]{Proposition}
\newtheorem{lemma}[theorem]{Lemma}
\newtheorem{corollary}[theorem]{Corollary}
\theoremstyle{remark}
\newtheorem{remark}[theorem]{Remark}
\DeclareMathOperator{\tr}{tr}
\newcommand{\R}{\mathbb R}
\newcommand{\cF}{\mathcal F}
\newcommand{\cZ}{\mathcal Z}
\newcommand{\dsharp}{d_{\dagger}}
\newcommand{\edag}{e_{\dagger}}
\newcommand{\adag}{a_{\dagger}}
\newcommand{\udag}{u_{\dagger}}
\newcommand{\doctype}{Preprint}
\newcommand{\docversion}{v0.01}
\newcommand{\docdate}{24 September 2026}
\newcommand{\versiondoi}{10.5281/zenodo.22940291}
\newcommand{\conceptdoi}{10.5281/zenodo.22940290}
\newcommand{\docrepo}{github.com/fsantibanezleal/CAOS\_RESEARCH}
\newcommand{\headerblock}{%
\begin{center}
\fbox{\parbox{0.94\textwidth}{\small\raggedright
\textbf{\doctype} \quad$\cdot$\quad \docversion, \docdate \quad$\cdot$\quad Licence CC BY 4.0\\[2pt]
CAOS Research open-problems programme \quad$\cdot$\quad problem \texttt{riemann-hypothesis}\\[2pt]
Version DOI \href{https://doi.org/\versiondoi}{\texttt{\versiondoi}}\\[2pt]
Concept DOI (always latest) \href{https://doi.org/\conceptdoi}{\texttt{\conceptdoi}}\\[2pt]
Not peer reviewed. Source code, data and computational records: \href{https://\docrepo}{\texttt{\docrepo}}
}}
\end{center}
\vspace{4pt}}
\pagestyle{fancy}
\fancyhf{}
\fancyhead[L]{\small CAOS Research \doctype}
\fancyhead[R]{\small Sharp three-point kernel\textperiodcentered\ \docversion}
\fancyfoot[C]{\thepage}
\fancypagestyle{plain}{\fancyhf{}\fancyfoot[C]{\thepage}\renewcommand{\headrulewidth}{0pt}}
\allowdisplaybreaks[1]

\title{A sharp three-point kernel bound and\\improved proportions of zeta zeros}
\author{Felipe Santib\'a\~nez-Leal\,\orcidlink{0000-0002-0150-3246}\\
\small CAOS Research program\\
\small\texttt{github.com/fsantibanezleal/CAOS\_RESEARCH}}
\date{\docdate}

\begin{document}
\maketitle
\thispagestyle{plain}
\headerblock

\begin{abstract}
Wang recently refined Lamzouri's Hilbert-space method by retaining a convex
spectral defect and proved that at least
$0.6725007703418699\ldots$ of the nontrivial zeros of the Riemann zeta
function are simple and on the critical line.  The quantitative input is a
three-point lower bound for a limiting Fourier kernel.  We determine the sharp
constant in the elementary ratio that controls this bound.  For all
$\alpha,\beta\ge0$,
\[
 \frac{\sqrt{(1+\alpha^2)(1+\beta^2)}
 +\alpha\sqrt{1+\alpha^2}+\beta\sqrt{1+\beta^2}}
 {1+\alpha^2+\alpha\beta+\beta^2}\le\sqrt2,
\]
with equality only at $(0,1)$ and $(1,0)$.  Substitution into Wang's argument
gives the simple-critical proportion
\[
 0.6725007995946757558283550562963947\ldots
\]
and the distinct-zero proportion
$0.8362503997973378779141775281481973\ldots$.  The new correction above the
Montgomery--Taylor constant is about $43.88\%$ larger than Wang's correction.
We also give a short-interval corollary for the source-certified rank-six
parity curve.  All numerical comparisons use directed rational enclosures and
an independent 120-decimal interval replay.  The global transfer depends on
Wang's recent unreviewed preprint.  No effective height, density-one result,
universal simplicity statement, or proof of the Riemann hypothesis is claimed.
\end{abstract}

\section{Introduction}

Let $N(T)$ count the nontrivial zeros $\rho=\beta+i\gamma$ of the Riemann
zeta function with $0<\gamma\le T$, including multiplicity.  Let $N_0^s(T)$
count the simple zeros among them on $\beta=1/2$, and let $N_d(T)$ count
distinct zeros, each once.  Define
\begin{equation}\label{eq:C0}
 C_0=\frac32-\frac1{\sqrt2}\cot(1/\sqrt2)
 =0.6725007036794116457343797908032951\ldots .
\end{equation}

Alp\"oge and Furman~\cite{AF}, through a proof whose central optimization was
found computationally, established that the lower density of simple critical
zeros is at least $C_0$.  Lamzouri~\cite{Lamzouri} gave a direct Hilbert-space
proof.  Wang~\cite{WangGlobal} retained a nonnegative spectral defect in
Lamzouri's finite inequality and combined it with disjoint three-point blocks.
His result is
\[
 \liminf_{T\to\infty}\frac{N_0^s(T)}{N(T)}
 \ge C_0+6.66624\ldots\times10^{-8}.
\]

The only change made here to Wang's analytic framework is an exact
optimization inside his three-point kernel lemma.  Wang obtains an auxiliary
ratio and bounds it above by $2$.  Theorem~\ref{thm:ratio} shows that its exact
maximum is $\sqrt2$.  This raises the kernel energy, and all subsequent steps
of the global proof remain valid with the new constant.

Our main consequence is the following.

\begin{theorem}\label{thm:global}
Assuming the analytic and finite-dimensional framework of
Wang~\cite{WangGlobal},
\begin{align}
 \liminf_{T\to\infty}\frac{N_0^s(T)}{N(T)}
 &\ge
 0.6725007995946757558283550562963947\ldots,
 \label{eq:global-simple}\\
 \liminf_{T\to\infty}\frac{N_d(T)}{N(T)}
 &\ge
 0.8362503997973378779141775281481973\ldots .
 \label{eq:global-distinct}
\end{align}
The displayed decimals are certified lower bounds after downward rounding.
\end{theorem}

The dependency phrase in Theorem~\ref{thm:global} is deliberate.  The sharp
ratio theorem below is self-contained, while the passage to zeta zeros uses a
September 2026 v1 preprint that has not been peer reviewed.  We inspect the
substitution point and reproduce the numerical constants, but do not relabel
the imported framework as independently established literature.

\section{The sharp auxiliary ratio}

For $\alpha,\beta\ge0$, set
\begin{equation}\label{eq:ratio}
 R(\alpha,\beta)=
 \frac{\sqrt{(1+\alpha^2)(1+\beta^2)}
 +\alpha\sqrt{1+\alpha^2}+\beta\sqrt{1+\beta^2}}
 {1+\alpha^2+\alpha\beta+\beta^2}.
\end{equation}

\begin{theorem}[Sharp ratio theorem]\label{thm:ratio}
For every $\alpha,\beta\ge0$,
\[
 R(\alpha,\beta)\le\sqrt2.
\]
Equality holds exactly at $(\alpha,\beta)=(0,1)$ and $(1,0)$.
\end{theorem}

\begin{proof}
Write
\[
 \alpha=\sinh u,\qquad \beta=\sinh v,
 \qquad u,v\ge0,
\]
and put
\[
 X=\cosh(u+v),\qquad Y=\cosh(u-v),\qquad
 S=\sinh(u+v)=\sqrt{X^2-1}.
\]
Then $X\ge Y\ge1$.  Product-to-sum identities give
\[
 \cosh u\cosh v=\frac{X+Y}{2},
 \qquad
 \sinh u\cosh u+\sinh v\cosh v=SY.
\]
Hence the numerator of~\eqref{eq:ratio} is
\begin{equation}\label{eq:num}
 \frac X2+Y\left(S+\frac12\right).
\end{equation}
A direct expansion gives the denominator as
\begin{equation}\label{eq:den}
 1+\sinh^2u+\sinh u\sinh v+\sinh^2v
 =\frac X2+Y\left(X-\frac12\right).
\end{equation}

Fix $X>1$.  The derivative of the quotient of~\eqref{eq:num} and
\eqref{eq:den} with respect to $Y$ has the sign of
\[
 1+S-X.
\]
This is positive, since $S>X-1$ follows by squaring from
$X^2-1>(X-1)^2$.  The quotient is therefore largest at $Y=X$.
Equality $Y=X$ means $\cosh(u-v)=\cosh(u+v)$, and for $u,v\ge0$ this is
equivalent to $u=0$ or $v=0$.

On this boundary the ratio is
\[
 \frac{1+\sqrt{X^2-1}}{X}.
\]
All quantities are nonnegative, and comparison with $\sqrt2$ reduces, after
two legitimate squarings, to
\[
 4(X^2-1)\le X^4
 \quad\Longleftrightarrow\quad
 (X^2-2)^2\ge0.
\]
Equality requires $X=\sqrt2$.  Together with $u=0$ or $v=0$, this gives the
two stated points.  If $X=1$, then $u=v=0$ and $R=1$.
\end{proof}

\begin{remark}
The parametrization is onto the whole nonnegative quadrant and the argument
includes all unbounded directions.  The proof is exact; no grid,
compactification, or numerical optimizer is used.
\end{remark}

\section{The improved three-point kernel bound}

We recall the exact point at which Theorem~\ref{thm:ratio} enters Wang's
argument.  Let $f_0$ be the limiting compactly supported cosine profile in
\cite{Lamzouri,WangGlobal}, let $K_0=\widehat f_0$, and define
\[
 a=\sqrt2\pi\cot(1/\sqrt2),
 \qquad \cF(x)=ax\sin(\pi x)-\cos(\pi x).
\]
The Fourier transform has the identity
\begin{equation}\label{eq:FK}
 \cF(x)=(2\pi^2x^2-1)K_0(x).
\end{equation}

For $u,v\ge0$, Wang sets
\[
 d=\max\{|\cF(u)|,|\cF(v)|,|\cF(u+v)|\},
 \qquad \alpha=au,\quad\beta=av.
\]
His vector triangle estimates and trigonometric addition identity yield
\begin{equation}\label{eq:wang-ratio-step}
 (1-d)^2\le d(1+d)R(\alpha,\beta)
\end{equation}
when $d<1$; the desired lower bound is automatic when $d\ge1$.
Wang then uses $R\le2$.  Theorem~\ref{thm:ratio} instead gives
\[
 (1-d)^2\le\sqrt2\,d(1+d).
\]
The positive equality root is
\begin{equation}\label{eq:ddag}
 \dsharp=
 \frac{\sqrt{2+8\sqrt2}-(2+\sqrt2)}{2(\sqrt2-1)}
 =0.2831654308085373270052981480781329\ldots .
\end{equation}
Thus $d\ge\dsharp$.  Exact directed enclosures also prove
\[
 \dsharp>\frac{\sqrt{57}-7}{2}>\sqrt5-2,
\]
so the result strictly improves both our initially declared fallback and
Wang's constant.

Combining~\eqref{eq:FK} with $u+v\le H$ gives the sharpened kernel lemma.

\begin{corollary}\label{cor:kernel}
For every $H>0$ and $u,v\ge0$ with $u+v\le H$,
\begin{equation}\label{eq:energy}
 K_0(u)^2+K_0(v)^2+K_0(u+v)^2
 \ge \edag(H):=
 \frac{\dsharp^2}{(1+2\pi^2H^2)^2}.
\end{equation}
\end{corollary}

\begin{proof}
At least one of the three $\cF$ values has modulus at least $\dsharp$.  For its
argument $x$, $|x|\le H$ and
$|2\pi^2x^2-1|\le1+2\pi^2H^2$.  Equation~\eqref{eq:FK} therefore makes the
corresponding squared kernel value at least the right side of
\eqref{eq:energy}.
\end{proof}

\section{Global proportions}

We give enough of the transfer to identify every dependence and constant.
For a conjugation-invariant finite zero multiset $\cZ$, let $n$ denote the
number of simple real elements and let
\[
 G_K=(K(x_j-x_\ell))_{1\le j,\ell\le n}
\]
be their Gram matrix.  Wang defines a nonnegative convex spectral defect
\[
 \Delta_K(\cZ)=\tr\Psi(G_K),
 \qquad
 \Psi(t)=
 \begin{cases}
 (t-1)^2,&0\le t\le2,\\
 2t-3,&t\ge2.
 \end{cases}
\]
His finite inequality is
\begin{equation}\label{eq:finite}
 n\ge2|\cZ|-\sum_{z,s\in\cZ}K(z-s)^2+\Delta_K(\cZ),
\end{equation}
with multiplicity in the sums.  A companion inequality gives the distinct
count.  His block lemma says that each disjoint simple-real triple whose
three squared pair kernels sum to at least $e$ contributes at least $2e$ to
the defect.

Partition the scaled ordinate interval into cells of length at most $H$ and
form disjoint triples in each cell.  At most two simple points per cell are
unused.  Corollary~\ref{cor:kernel} and the Riemann--von Mangoldt asymptotic
then give, after the same smoothing and ordered limits as in
\cite{WangGlobal},
\begin{equation}\label{eq:globalformula}
 \liminf_{T\to\infty}\frac{N_0^s(T)}{N(T)}
 \ge\udag(H):=
 \frac{C_0-2\adag(H)/H}{1-\adag(H)},
 \qquad
 \adag(H)=\frac{2\edag(H)}3.
\end{equation}
The smoothing parameter remains fixed while $T$ tends to infinity and is
removed only afterward.  Since~\eqref{eq:energy} is uniform on every fixed
cell, the sharper constant does not alter that order of limits.

Choose the rational value
\begin{equation}\label{eq:Hstar}
 H_*=\frac{372019}{100000}=3.72019.
\end{equation}
Directed arithmetic gives
\begin{align*}
 \udag(H_*)-C_0
 &=0.0000000959152641100939752654930995979\ldots,\\
 \udag(H_*)
 &=0.6725007995946757558283550562963947865\ldots .
\end{align*}
This proves~\eqref{eq:global-simple}.

For comparison, using Wang's $\sqrt5-2$ and his exact
$H_0=372018941724/10^{11}$ in the same certificate gives
\[
 \delta_0=
 0.0000000666624583334125249521207539241\ldots,
\]
strictly between $6.66624\times10^{-8}$ and
$6.66625\times10^{-8}$.  The lower endpoint of the new gain interval exceeds
the upper endpoint of this reproduced value.  Their ratio shows that the new
correction is about $43.88\%$ larger.  We make no claim that the convenient
rational $H_*$ is the exact optimizer.

Wang's distinct-zero inequality, with the same replacement, is
\[
 \liminf_{T\to\infty}\frac{N_d(T)}{N(T)}
 \ge\frac{1+\udag(H)}2.
\]
Inserting~\eqref{eq:Hstar} proves~\eqref{eq:global-distinct}.

\section{A short-interval corollary}

The same cell packing can be placed inside the fixed-power intervals used in
the companion CAOS manuscript~\cite{CAOSShort}.  Let $h_6(\theta)$ be its
rank-six Hilbert--parity lower curve, $k_6(\theta)$ the attributed odd-support
term, and $q(\theta)$ the pair statistic.  If $h_6(\theta)>0$, choose a fixed
cell length $H>2/h_6(\theta)$ and set
\[
 \alpha_H=\frac{2\edag(H)}3,
 \qquad
 \beta_H=\frac{2\alpha_H}{H}.
\]
The number of cells divided by the local zero count tends to $1/H$.  The
triple count and Wang's block lemma therefore give
\[
 \liminf\frac{\Delta_K}{N}\ge\alpha_Hs-\beta_H,
\]
where $s$ is the lower simple-critical proportion.  Substitution into the
finite defect--parity product of~\cite{CAOSShort} yields
\begin{equation}\label{eq:shortquad}
 2(1-s)^2\le(1-k_6(\theta))
 \{q(\theta)+\beta_H-(1+\alpha_H)s\}.
\end{equation}

At $\theta=0.5459$ and the integer cell length $H=140730$, the exact
certificate proves $H>2/h_6$.  The smaller root $J_6$ of equality in
\eqref{eq:shortquad} satisfies
\begin{equation}\label{eq:shortgain}
 J_6-h_6>
 3.08678099833341875132438420197419\times10^{-31}.
\end{equation}
The bound is correlated: it follows from the value of the strengthened
quadratic at the old root and a derivative bound, rather than subtraction of
two independently rounded roots.  It exceeds the earlier spectral correction
$1.776662254114\ldots\times10^{-68}$ by more than 37 orders of magnitude.
It remains too small to alter the useful decimals of $h_6$ and does not move
the positivity onset.

\section{Exact certificate and limitations}

The machine-readable record~\cite{RecordIX} pins both files of
arXiv:2609.24167v1, the experiment declaration and amendment, the preceding
rank-six certificate, the runner, and the canonical output.  The principal
calculation uses exact fractions, integer square-root enclosures, and directed
Taylor bounds for elementary functions.  A separate implementation using
120-decimal \texttt{mpmath.iv} intervals overlaps every headline interval.
The canonical CPU run completed in 37.157 seconds from a clean Git commit.

The certificate verifies
\begin{itemize}
\item the exact ordering of all three kernel constants;
\item Wang's printed correction at his rational cell length;
\item strict interval separation between the new and old global gains;
\item the short-interval cell condition and positive correlated reserve;
\item strict dominance over the preceding spectral correction; and
\item all pinned source and execution hashes.
\end{itemize}

Theorem~\ref{thm:ratio} is unconditional.  The global transfer is conditional
on the correctness of the cited Wang v1 argument, although the exact
substitution point and downstream algebra were rechecked.  The short-interval
corollary additionally imports a rank-six constant whose source does not print
its coefficient vector.  The statements are asymptotic and supply no
effective starting height.  They neither show that every zero lies on the
critical line nor that every zero is simple.

\section*{Computational assistance disclosure}

ChatGPT-6 Astra was used to search the recent literature, explore the
hyperbolic substitution, implement interval certificates, and check the
manuscript.  The author inspected the primary sources, verified the proof and
equality cases, reviewed every stated dependency and limitation, and takes
responsibility for the content.  No artificial-intelligence system is listed
as an author.

\begin{thebibliography}{99}
\raggedright

\bibitem{AF}
L. Alp\"oge and R. Furman,
\emph{More than two thirds of the zeta zeros are simple and on the critical line},
arXiv:2608.13637v2 (2026),
\url{https://arxiv.org/abs/2608.13637v2}.

\bibitem{CAOSShort}
F. Santib\'a\~nez-Leal,
\emph{Simple critical zeros in short intervals: stability, parity,
localization, Hilbert compression, and spectral defect}, v0.07,
Zenodo (2026),
\url{https://doi.org/10.5281/zenodo.22860012}.

\bibitem{Lamzouri}
Y. Lamzouri,
\emph{A new proof that more than $2/3$ of the zeros of the Riemann zeta
function are simple and on the critical line},
arXiv:2609.02882v2 (2026),
\url{https://arxiv.org/abs/2609.02882v2}.

\bibitem{RecordIX}
F. Santib\'a\~nez-Leal,
\emph{EXP-009: Wang kernel sharpening and parity transfer},
CAOS Research computational record (2026),
\url{https://github.com/fsantibanezleal/CAOS_RESEARCH/tree/main/problems/number-theory/riemann-hypothesis/experiments/EXP-009-wang-kernel-sharpening}.

\bibitem{Titchmarsh}
E. C. Titchmarsh,
\emph{The Theory of the Riemann Zeta-function}, second edition,
revised by D. R. Heath-Brown, Oxford University Press, 1986.

\bibitem{WangGlobal}
B. Wang,
\emph{Proportions of the non-trivial zeros of the Riemann zeta function},
arXiv:2609.24167v1 (2026),
\url{https://arxiv.org/abs/2609.24167v1}.

\bibitem{WangShort}
B. Wang,
\emph{Simple critical zeros and distinct zeros of the Riemann zeta-function
in short intervals},
arXiv:2609.07918v1 (2026),
\url{https://arxiv.org/abs/2609.07918v1}.

\end{thebibliography}
\end{document}
